% ==============================================================================
% Abstract
% ==============================================================================
\chapter*{Abstract}
\label{ch:abstract}

Street vending sustains informal urban livelihoods across developing economies while supplying affordable goods and essential sustenance to millions of city residents. In Indian metropolitan regions, statutory mandates such as the \textit{Street Vendors (Protection of Livelihood and Regulation of Street Vending) Act, 2014} require municipal bodies to legally demarcate vending zones, prevent eviction of registered vendors, and balance public space allocation. However, municipal town-planning authorities lack objective, data-driven decision-support frameworks to evaluate zone suitability, enforce statutory carrying capacity, execute equitable relocations, and engineer physical micro-scale stall spaces. Uncoordinated street vending consequently results in pedestrian walkway encroachments, vehicular congestion, and recurring civic disputes.

This report presents the design, full-stack implementation, and empirical validation of \textbf{VendorMap} (UrbanVending AI)---a smart vendor relocation and zone recommendation platform developed in collaboration with municipal realities in Nashik, Maharashtra. Evaluated on official civic datasets from the Nashik Municipal Corporation (NMC) encompassing 7,259 registered vendors and 308 municipal zones across six administrative divisions, the platform combines OpenStreetMap (OSM) GIS spatial feature extraction (13,988 road links, 353 commercial POIs) with unsupervised $K$-Means clustering ($K=2$, Adjusted Rand Index $ARI = 1.0$) to establish baseline environmental typologies without confusing unsupervised patterns with legal zoning rules.

To compute robust relocation recommendations, an operational six-factor Multi-Criteria Decision Analysis (MCDA) engine combines logarithmic saturation transforms, exponential proximity decay, and empirical Bayes category shrinkage ($M=20.0$) with leave-location-group-out cross-validation, calibrated via the Analytic Hierarchy Process (AHP) ($CR \le 0.10$). In a 224-scenario sensitivity audit, the MCDA engine achieved 100\% Top-1 stability and a mean Spearman rank correlation of $\rho = 0.9977$. To resolve the \textit{Spatial Capacity Paradox}---wherein greedy unconstrained suitability recommendation causes catastrophic 900\% capacity over-allocation in prime commercial cores---the system integrates two-stage pre-emptive Goal Programming (MILP via HiGHS), guaranteeing zero capacity violation while minimizing unassigned vendors and preferred-division fallback penalties.

Physical feasibility is demonstrated through a metric micro-slot packing generator validated on the Lekha Nagar pilot site (Zone NN\_F\_013), successfully packing 10 standardized $2\,\text{m} \times 2\,\text{m}$ stall footprints with sub-millimeter precision ($1.87 \times 10^{-9}\,\text{m}$ edge error), $0.5\,\text{m}$ inter-slot spacing, and a continuous $1.2\,\text{m}$ pedestrian clearance corridor in UTM Zone 43N metric projection. Supervised benchmarking on 5-fold grouped cross-validation confirms that Random Forest regression ($R^2 = 0.7592$) achieves a 59.8\% error reduction over baselines, with TreeSHAP explainability identifying proximate commercial footfall as the dominant driver of suitability.

The operational delivery system features an authenticated React 18 / FastAPI service architecture with SQLite Write-Ahead Logging, strict cryptographic data minimization (HMAC-SHA256 identity hashing for vendor mobile numbers), in-app turn-by-turn walking navigation using Leaflet and GPS geolocation, and an official, direct integration with the Meta WhatsApp Cloud API for bilingual (English/Marathi) conversational access. The result is an institute-compliant, auditable, and privacy-preserving platform for smart municipal governance and street vendor livelihood protection.

\vspace{0.8cm}
\noindent \textbf{Keywords:} Street Vendor Relocation, Geospatial Decision Support, Multi-Criteria Decision Analysis (MCDA), Analytic Hierarchy Process (AHP), Mixed-Integer Linear Programming (MILP), Goal Programming, Micro-Slot Demarcation, Geographic Information Systems (GIS), Explainable AI (XAI), TreeSHAP, WhatsApp Cloud API.
